%% This section adds chapter 1
\section{Introduction}
% example entries for nomenclature. These entries should be given throughout the text right after the symbol is used the first time. If "refpage" option is given, the corresponding page number is shown in nomenclature. The following nomenclature groups are used; A=Latin, B=Greek, C=Subscripts, D=Superscripts, E=Abbreviations, F=Dimensionless. This letter code should be placed first inside the [] brackets. The following letters inside square brackets [] are used to alphabetically sort the symbols and abbreviations in the list. The symbol or abbreviation itself is defined within the curly braces {} after the square brackets. The $$ signs create an equation form, i.e. a variable in italics. Inside the latter brackets, the legend and the unit are defined with the \nounit{} definition. For units, italics can be removed with the \mathrm{} command. See instructions https://www.overleaf.com/learn/latex/Nomenclatures or https://www.ctan.org/pkg/nomencl. 
% A=Latin letters
\nomenclature[Aa1]{$A$}{area\nomunit{$\mathrm{m^2}$}}
\nomenclature[Aa2]{$a$}{constant}
\nomenclature[Acd]{$C_\mathrm{D}$}{drag coefficient}
\nomenclature[Acp]{$c_p$}{specific heat capacity at constant pressure\nomunit{$\mathrm{J/(kgK)}$}}
\nomenclature[Acv]{$c_v$}{specific heat capacity at constant volume\nomunit{$\mathrm{J/(kgK)}$}}
\nomenclature[Ad]{$d$}{diameter\nomunit{$\mathrm{m}$}}
\nomenclature[AF]{$\mathbf{F}$}{force vector\nomunit{$\mathrm{N}$}}
\nomenclature[Af]{$f$}{frequency\nomunit{$\mathrm{1/s,Hz}$}}
\nomenclature[Ag]{$g$}{acceleration due to gravity\nomunit{$\mathrm{m/s^2}$}}
\nomenclature[Ah]{$h$}{heat transfer coefficient\nomunit{$\mathrm{W/(m^2K)}$}}
%\nomenclature[Ah2]{$h$}{specific enthalpy\nomunit{$\mathrm{J/kg}$}}
\nomenclature[Aj]{$\mathbf{j}$}{flux vector\nomunit{$\mathrm{m/s}$}}
\nomenclature[AL]{$L$}{characteristic length\nomunit{$\mathrm{m}$}}
\nomenclature[Al]{$l$}{length\nomunit{$\mathrm{m}$}}
\nomenclature[AM]{$M$}{torque\nomunit{$\mathrm{Nm}$}}
\nomenclature[Am]{$m$}{mass\nomunit{$\mathrm{kg}$}}
\nomenclature[AN]{$N$}{number of particles}
\nomenclature[An2]{$\mathbf{n}$}{unit normal vector}
\nomenclature[Ap]{$p$}{pressure\nomunit{$\mathrm{Pa}$}}
\nomenclature[Aq]{$q$}{heat flux\nomunit{$\mathrm{W/m^2}$}}
\nomenclature[Aqm]{$q_\mathrm{m}$}{mass flow rate\nomunit{$\mathrm{kg/s}$}}
\nomenclature[Ar]{$r$}{radius\nomunit{$\mathrm{m}$}}
\nomenclature[AT]{$T$}{temperature\nomunit{$\mathrm{K}$}}
\nomenclature[At]{$t$}{time\nomunit{$\mathrm{s}$}}
\nomenclature[AV]{$V$}{volume\nomunit{$\mathrm{m^3}$}}
\nomenclature[Av]{$v$}{velocity magnitude\nomunit{$\mathrm{m/s}$}}
\nomenclature[Av2]{$\mathbf{v}$}{velocity vector\nomunit{$\mathrm{m/s}$}}
\nomenclature[Ax]{$x$}{x-coordinate (width)\nomunit{$\mathrm{m}$}}
\nomenclature[Ay]{$y$}{y-coordinate (depth)\nomunit{$\mathrm{m}$}}
\nomenclature[Az]{$z$}{z-coordinate (height)\nomunit{$\mathrm{m}$}}
% B=Greek symbols
%\nomenclature[Ba]{$\alpha$}{thermal expansion coefficient\nomunit{$\mathrm{1/K}$}}
\nomenclature[Ba]{$\alpha$}{(alfa)\nomunit{$\mathrm{1/K}$}}
\nomenclature[Bb]{$\beta$}{(beta)}
\nomenclature[BG]{$\Gamma$}{(Gamma)}
\nomenclature[Bg]{$\gamma$}{(gamma)}
\nomenclature[BD1]{$\Delta$}{(Delta) usually used for change/difference}
\nomenclature[BD2]{$\varDelta$}{(varDelta)}
\nomenclature[Bd3]{$\delta$}{(delta) notice the difference to $\partial$ (partial differential)}
\nomenclature[Be1]{$\epsilon$}{(epsilon)}
\nomenclature[Be2]{$\upepsilon$}{(upepsilon)}
\nomenclature[Bz]{$\zeta$}{(zeta)}
\nomenclature[Be3]{$\eta$}{(eta)}
\nomenclature[Bth]{$\theta$}{(theta)}
\nomenclature[BTh]{$\Theta$}{(Theta)}
\nomenclature[Bth2]{$\vartheta$}{(vartheta)}
\nomenclature[Bi]{$\iota$}{(iota)}
\nomenclature[Bk]{$\kappa$}{(kappa)}
\nomenclature[BL]{$\Lambda$}{(Lambda)}
\nomenclature[Bl]{$\lambda$}{(lambda)}
\nomenclature[Bm]{$\mu$}{(mu)}
\nomenclature[Bm2]{$\upmu$}{(upmu)}
\nomenclature[Bn]{$\nu$}{(nu)}
\nomenclature[BX]{$\Xi$}{(Xi)}
\nomenclature[Bx]{$\xi$}{(xi)}
\nomenclature[BP]{$\Pi$}{(Pi)}
\nomenclature[Bp]{$\pi$}{(pi) e.g. pressure ratio, note the difference with (uppi) constant $\uppi=3.14\ldots$}
\nomenclature[Br]{$\rho$}{(rho)}
\nomenclature[Br2]{$\uprho$}{(uprho)}
\nomenclature[BS]{$\Sigma$}{(Sigma)}
\nomenclature[Bs]{$\sigma$}{(sigma)}
\nomenclature[Bs2]{$\varsigma$}{(varsigma)}
\nomenclature[Bta]{$\tau$}{(tau)}
\nomenclature[Bu]{$\upsilon$}{(upsilon)}
\nomenclature[Bph]{$\phi$}{(phi)}
\nomenclature[BPh]{$\Phi$}{(Phi)}
\nomenclature[Bph2]{$\varphi$}{(varphi)}
\nomenclature[BX]{$\chi$}{(chi)}
\nomenclature[BPs]{$\Psi$}{(Psi)}
\nomenclature[Bps]{$\psi$}{(psi)}
\nomenclature[BO]{$\Omega$}{(Omega)}
\nomenclature[Bo]{$\omega$}{(omega)}
% C=Subscripts
\nomenclature[Ce]{eff}{effective}
\nomenclature[Cma]{max}{maximum}
\nomenclature[Cmi]{min}{minimum}
\nomenclature[Cg]{g}{gas}
\nomenclature[Ct]{tot}{total}
\nomenclature[Cp]{p}{particle}
\nomenclature[Cs]{s}{solid}
\nomenclature[Cl]{l}{liquid}
% D=Superscripts
\nomenclature[Dp]{p}{partial layer}
\nomenclature[Ds]{*}{dimensionless}
% E=Abbreviations
\nomenclature[E2]{2D}{two dimensional}
\nomenclature[E3]{3D}{three dimensional}
\nomenclature[Ec]{CFD}{computational fluid dynamics}
\nomenclature[El]{LES}{large eddy simulation}
\nomenclature[Ep]{PDF}{probability density function}
% F=Dimensionless
\nomenclature[FAr]{Ar}{Archimedes number}
\nomenclature[FBi]{Bi}{Biot number}
\nomenclature[FFO]{Fo}{Fourier number}
\nomenclature[FGr]{Gr}{Grashof number}
\nomenclature[FNu]{Nu}{Nusselt number}
\nomenclature[FPr]{Pr}{Prandtl number}
\nomenclature[FRe]{Re}{Reynolds number}
\nomenclature[FSh]{Sh}{Sherwood number}
\nomenclature[FSt1]{St}{Stanton number}
\nomenclature[FSte]{Ste}{Stefan number}
\nomenclature[FStk]{Stk}{Stokes number}
%%%%%%%%%%%%%%%%%%%%%%%%%%%%%%%%%
% Start of Introduction section %
%%%%%%%%%%%%%%%%%%%%%%%%%%%%%%%%%
\subsection{CBRN Defense Field}
CBRN defense field encompasses the military and civil measures taken to protect against chemical, biological, radiological, and nuclear threats. The field includes a wide range of activities, such as detection, identification, decontamination, and medical countermeasures \citep{Otrisal2014}. CBRN events can pose a significant physical and psychological hazard to the public and pose environmental and economic impact \citep{MOHMalaysia2025}. Fast, coordinated and accurate response in the first critical hour or two is crucial to minimize the impact of such events \citep{JESIP2016}. Providing accurate information in a timely manner is essential for effective decision-making and response \citep{MOHMalaysia2025}.
\par
Coordination between the various agencies and personnel involved in CBRN defense is essential. Numerous software solutions have been developed to facilitate communication and information sharing among CBRN defense teams \citep{Marques2025,EnvironicsEnviScreenX,BruhnNewTechFrontline,ObservisObsasLink}. A common pattern among these solutions is the near real-time consolidation and visualization of data from multiple sources, such as sensors, first responders, command centers and various governmental databases. However, the effectiveness of these solutions is  limited by the communication infrastructure available in the affected area. Additionally, any delays introduced by the communication infrastructure, protocols, and data processing can add up to significant delays in the overall response time. In the case of CBRN events, where every second counts, these delays can have serious consequences.
\subsection{LEO Satellite Networks}
One of the more recent developments in communication infrastructure is the commercial availability of low earth orbit (LEO) satellite internet. LEO satellite networks have the potential to provide high-speed, low-latency internet access to remote and underserved areas \citep{DuanTong2021SSNC}. This technology has the potential to improve the communication abilities of CBRN defense teams by providing a communication channel in remote areas or areas with delibierately disrupted communication infrastructure \citep{AbelsJoscha2024Piig}.
\par
One of the leading providers of LEO satellite internet is Starlink, which as of September 2024 has launched 6426 satellites into orbit and has 14 million active subscribers globally, making it the leader by a wide margin in both of these metrics. Starlink subscriptions are widely available across the globe in both personal and business plans \citep{KooshaBehzad2025CAoR}. The relatively low subscription cost, easy of acquisition, low-latency, high bandwidth and coverage make it a compelling option as the communication infrastructure for CBRN defense teams.
\subsection{QUIC Protocol}
QUIC is a multiplexed, encrypted and low-latency transport protocol that is designed with the goal of improving transport performance and security of internet applications. QUIC is a cross-layer protocol and operates at the Transport, Security and Application layers of the OSI model \citep{LangleyAdam2017TQTP}.
\par
According to the designers of the QUIC protocol, two of the main motivators behind its design were reducing the handshake delay of TCP and TLS and Head-of-line blocking delay within TCP. Additionally, TCP is a rigid protocol which cannot be modified to suit the needs of modern applications or specific use cases. QUIC comes out as an answer to these limitations and by measurements from Google it reduces latency of Google searches by 8\% on desktop and by 3.6\% on mobile. \citep{LangleyAdam2017TQTP} The flexibility and low latency characteristics of QUIC makes it particularly suitable for the demanding application of CBRN monitoring and communication over LEO satellite networks.
\subsection{Motivation and Context}
 This thesis aims to address the challenges of near real-time communication for CBRN applications with global coverage by proposing a computational design science approach for inter-vehicle network communication using QUIC as a transport protocol over LEO satellite internet. The proposed approach leverages the advantages of LEO satellite internet, such as low latency and high bandwidth, to enhance the communication capabilities of CBRN defense vehicles. QUIC is used as a transport protocol due to its flexibility for specific applications and low latency. By improving the speed and accuracy of information sharing, the proposed approach has the potential to significantly enhance the effectiveness of CBRN defense operations. The proposed design aims to enable wider capabilities of software systems for CBRN defense, such as global connectivity in remote areas and reduced latency and reliability over traditional network infrastructure (e.g. 3G, 4G, LTE, 5G and geostationary satellite internet) and transport protocol.

\subsection{Objectives and Scope of This Work}
This thesis adopts the Design Science Research Methodology (DSRM) framework by \citet{DSRMethodology2007} to design, implement, and evaluate a proof-of-concept (PoC) software communication artifact. The primary aim of the artifact is to leverage the QUIC transport protocol over Low Earth Orbit (LEO) satellite internet for near real-time CBRN communication. Furthermore, Machine Learning (ML) will be used to optimize the performance of the QUIC protocol over LEO satellite networks, with a focus on minimizing latency, improving resilience and prioritizing critical CBRN data payloads. The artifact will be evaluated in a real-world environment using Starlink's commercial LEO satellite network.

To achieve this aim, the work is divided into the following concrete research objectives:
\begin{itemize}
    \item \textbf{O1 (Artifact Design and Implementation):} Develop a software prototype capable of transmitting and receiving CBRN data payloads between remote vehicular nodes using the QUIC transport protocol with ML optimizations.
    \item \textbf{O2 (Payload Simulation):} Design a standardized, synthetic data schema that mimics real-world CBRN sensor telemetry (e.g., detection alerts, GPS coordinates, and timestamped hazard measurements) without relying on proprietary hardware.
    \item \textbf{O3 (Empirical Evaluation):} Measure the performance of the implemented QUIC artifact over a live commercial LEO satellite network (Starlink) under real-world operating conditions, tracking metrics such as end-to-end latency, handshake duration, packet loss, and connection recovery during satellite handovers.
    \item \textbf{O4 (Comparative Benchmarking):} Compare the performance of the QUIC-based artifact with and without ML optimizations against a traditional baseline protocol stack (i.e., TCP with TLS 1.3) running over the same LEO satellite network.
\end{itemize}

\par
To guide the empirical evaluation of the artifact, this study formulates two core research questions (RQs) and their corresponding testable hypotheses ($H_1$ and $H_2$):

\begin{quote}
    \textbf{RQ1:} How does replacing traditional TCP/TLS transport protocols with standard QUIC affect the latency, connection resilience, and real-time data delivery of CBRN telemetry when routed over commercial LEO satellite networks?
    
    \textbf{Hypothesis 1 ($H_1$):} Because of its 0-RTT/1-RTT handshake mechanism and elimination of head-of-line blocking, utilizing the standard QUIC protocol over LEO satellite internet significantly reduces end-to-end transmission latency and improves connection recovery times during network dropouts compared to a standard TCP/TLS baseline.
    
    \vspace{0.5cm}
    
    \textbf{RQ2:} To what extent does applying Machine Learning (ML) optimizations to the QUIC protocol further enhance packet prioritization and resilience during predictable LEO network degradations, such as satellite handovers?
    
    \textbf{Hypothesis 2 ($H_2$):} An ML-optimized QUIC implementation will yield a statistically significant reduction in latency for high-priority CBRN payloads during satellite handovers compared to a non-optimized QUIC baseline, by dynamically adapting transport parameters to the fluctuating LEO network state.
\end{quote}

\par
\textbf{Delimitations and Scope:}
To maintain scientific rigor and feasibility within the timeframe of a Master's thesis, explicit boundaries are set regarding the scope of this study:
\begin{itemize}
    \item \textbf{Network Infrastructure:} The empirical testing is strictly limited to the Starlink commercial LEO satellite constellation. Other LEO constellations (such as OneWeb or Amazon Kuiper), geostationary (GEO) satellites, and terrestrial cellular networks (e.g. 3G, 4G, LTE, 5G) are excluded from the experimental setup.
    \item \textbf{Hardware and Sensors:} Physical CBRN detection hardware, chemical analyzers, and vehicular sensor integration are out of scope. The study relies entirely on software-simulated data payloads that accurately replicate the data volume and frequency of standard CBRN defense devices.
    \item \textbf{C4I Integration and UI:} While the artifact demonstrates data transport capabilities, developing a full Command, Control, Communications, Computers, and Intelligence (C4I) situational awareness dashboard or graphical user interface (GUI) is out of scope. The focus remains strictly on the networking and transport layers.
    \item \textbf{Adversarial Security Testing:} Although QUIC inherently enforces TLS 1.3 encryption, active electronic warfare testing—such as intentional radio frequency (RF) jamming, satellite spoofing, or deep cryptographic penetration testing—is excluded from this work.
\end{itemize}

\clearpage  % This command will start the next section from the new page
